%% [migrated] heading absorbed into chapter wrapper: Basic algorithms for LQFT
A good reference for more detail is [L\"uscher 1002.4232]
\begin{itemize}
  \item The basic numerical task in LQFT is evaluation of correlation functions which corresponds to very high dimensional integrals with specific structure
\end{itemize}

\begin{equation}
\langle\mathcal{O}(\phi)\rangle=\frac{1}{Z} \int \mathcal{D} \phi e^{-S[\phi]} \mathcal{O}[\phi]
\end{equation}

where only $\mathcal{O}$ changes for different physical quantities in the same theory

\begin{itemize}
  \item For fermionic theories, Grassman integration is done analytically and requires determinants and inverses of the Dirac operator
\end{itemize}

\begin{align}
    \langle \mathcal{O}(\bar{\psi}, \psi, \phi) \rangle &= \frac{1}{Z} \int \mathcal{D}\phi \mathcal{D}\bar{\psi} \mathcal{D}\psi e^{-S[\phi]-\int \bar{\psi}D[\phi] \psi} \mathcal{O}[\phi, \bar{\psi}, \psi] \nonumber \\
    &= \frac{1}{Z}\int \mathcal{D}\phi  \operatorname{det}[D[\phi]]e^{-S[\phi]}\mathcal{O}[\phi, D^{-1}[\phi]] \nonumber \\
    &= \frac{1}{Z}\int \mathcal{D}\phi  e^{-S_\text{eff}[\phi]}\mathcal{O}[\phi],
\end{align}
where $S_\text{eff}[\phi] \equiv S[\phi] + \ln \operatorname{tr}[D[\phi]]$.

\begin{itemize}
    \item For a typical lattice geometry that is used of $100^4$ with $N_c = 3$, the integration is over $3 \times 10^9$ degrees of freedom and clearly cannot be done using quadrature
\end{itemize}

\section{Importance Sampling}

\begin{itemize}
  \item Given a simple $1$D integral
\end{itemize}

\begin{equation}
I=\int_{a}^{b} \text{d} x f(x)
\end{equation}

if \(f(x)\) is smooth this is easily approximated by quadrature

\begin{equation}
\bar{I}=\frac{b-a}{N} \sum_{i=1}^{N} f\left(x_{i}\right)=(b-a)\langle f\rangle
\end{equation}

where \(x_{i}\) are uniformly distributed over \([a, b]\), e.g. \(x_{i}=a+\frac{b-a}{N}(i-\frac{1}{2})\) with

\begin{equation}
\bar{I}=I+\mathcal{O}\left(\frac{\sigma}{\sqrt{N}}\right) \quad \text { with } \sigma^{2}=\frac{1}{N-1} \sum_{i=1}^{N}\left(f\left(x_i\right)-\langle f\rangle\right)^{2}
\end{equation}

However the value of \(N\) for which convergence become asymptotic depends strongly on the structure of \(f\) (and for some $f$'s is not guaranteed)

\begin{itemize}
  \item For any \(f(x)\) we are free to change variables to improve the convergence of numerical quadrature. If
\end{itemize}

\begin{equation}
f(x)=g(x) p(x)
\end{equation}

where \(p(x)\) is positive and bounded with \(\int_{a}^{b} p(x) \text{d} x=1\) then we can change variables to

\begin{equation}
y=\int_{a}^{x} p(z) \text{d} z
\end{equation}

and define \(\tilde{g}(y)=g(x(y))\) so that
\begin{equation}
I=\int_{a}^{b} f(x) \text{d} x=\int_{a}^{b} g(x) p(x) \text{d} x=\int_{0}^{1} \tilde{g}(y)  \text{d} y \simeq \frac{1}{N} \sum_{i=1}^{N} \tilde{g}\left(y_{i}\right)
\end{equation}

where \(y_{i}\) are uniformly distributed on \([0,1]\). NB: an integrated form here is used since if $\hat{y} = p(x)$ was used and $p(x) = 0$ somewhere then $x(y)$ does not exist. If $p \neq 0$ anywhere then this can be done.

With this change, the uniform sampling in $y$ corresponds to sampling in $x$ that is concentrated in regions of high probability and thus this sampling will converge more quickly.

\begin{itemize}
    \item In QFT the integrals are high dimensional versions of this scenario with the probability measure
\end{itemize}
\begin{equation}
    \text{d}P[\phi] = \begin{cases}
        e^{-S[\phi]}\mathcal{D}\phi \qquad \qquad \: \: \:\text{purely bosonic}\\
        e^{-S[\phi]}\operatorname{det}D[\phi]\mathcal{D}\phi \qquad \text{including fermions}
    \end{cases}
\end{equation}


These probability measures correlate the variables $\phi(x)$ and inverting $x = P(\phi)$ is not possible. However if we could generate samples of $\phi$ with the distribution $P(\phi)$ then we could use importance sampling

\section{Markov Chains}
\begin{itemize}
    \item A Markov chain is a stochastic sequence of states/configurations of the degrees of freedom of a system defined by a transition probability $T(s,s')$ from a state $s'$ to a state $s$ in which the transition only depends on $s$ and $s'$ and not on other states in the chain:
\end{itemize}
\begin{equation}
    T: \mathcal{S} \rightarrow \mathcal{S}, \quad \text{where }\mathcal{S}\text{ is the space of all possible states}
\end{equation}
\begin{figure}
    \centering
    \redrawcompare{fig-notes-078.png}{fig-redraw-078.pdf}
    \caption{A Markov chain (red) containing states (denoted by points) in a state space (green).}
    \label{fig:fig_mc}
\end{figure}
The transition probabilities satisfy
\begin{equation}
\begin{aligned}
& \text { 1) } \quad \sum_{s^{\prime}} T\left(s^{\prime}, s\right)=1 \quad \forall s\\
& \text { 2) } \quad 0 \leqslant T\left(s^{\prime}, s\right) \leqslant 1 \quad \forall s^{\prime}, s
\end{aligned}
\end{equation}

\begin{itemize}
  \item An ensemble is an infinite set of sample states \(E=\left\{s_{1}, s_{2}, \ldots\right\}\) with a density given by a probability distribution \(P_{E}(s)\)
  \begin{itemize}
      \item Given an ensemble, $E$, the transition probabilities are such that 
      \begin{equation}
          E \underset{T}{\rightarrow} E^\prime=\{s'=T(s',\ )P_E(s) \forall s \in E \}
      \end{equation} 
      which has density $P_{E'}(s)$. The transition probability is also a map $P_E \underset{T}{\rightarrow}P_{E'}$
  \end{itemize}
\end{itemize}

\begin{itemize}
  \item To solve our sampling problem, we want to generate a set of samples distributed according to an equilibrium probability distribution $P_\text{eq}(s)$ that is given by the Boltzmann factor $e^{-S(\phi)}$ with a state $s$ representing the degrees of freedom of the fields $\phi(x) \:\forall x \in \Lambda$
  \item To generate these samples, a Markov process needs to have two properties
\end{itemize}
\begin{enumerate}
  \item The process must have \(P_{\text{eq}}\) as a fixed point. That is
\end{enumerate}

\begin{equation}
\sum_{s} P_{\text{eq}}(s) T\left(s^{\prime}, s\right)=P_{\text{eq}}\left(s^{\prime}\right)
\end{equation}

so that once the chain has reached this distribution, it does not leave it

\begin{itemize}
  \item Since \(\sum_{s} T\left(s, s^{\prime}\right)=1\) this implies the balance condition
\end{itemize}

\begin{equation}
\sum_{s} P_{\text{eq}}(s) T\left(s^{\prime}, s\right)=\sum_{s} T\left(s, s^{\prime}\right) P_{\text{eq}}\left(s^{\prime}\right),
\end{equation}
as can be seen by multiplying the RHS by $1$; $s'$ is fixed. This says that the probability of reaching a state equals the probability of leaving it.

\begin{itemize}
  \item The balance condition is guaranteed by the detailed balance condition (easier to impose)
\end{itemize}

\begin{equation}
T(s', s) P_{\text{eq}}(s)=T\left(s, s^{\prime}\right) P_{\text{eq}}\left(s^{\prime}\right)
\end{equation}

\begin{enumerate}
  \setcounter{enumi}{1}
  \item Since the integral that is being sampled is over all values of the fields the Markov process should be ergodic: it is possible to get to every state. Thus $0 < T(s, s')^n < 1 \quad \forall s,s'$ and some $n$. It does not need to be a single step to get to every state.
\end{enumerate}
\begin{itemize}
    \item A Markov process satisfying these criteria is guaranteed to reach $P_\text{eq}$ after many steps starting from any state. After reaching $P_\text{eq}$, new states are produced in the sequence with this probability distribution
    \begin{itemize}
        \item To see this, define the distance between ensembles through their distributions
    \end{itemize}
\end{itemize}

\begin{equation}
\left\|E_{a}-E_{b}\right\| \equiv \sum_{s}\left|P_{a}(s)-P_{b}(s)\right|
\end{equation}

Then given an ensemble \(E\), a Markov step acting on \(E\) produces a new ensemble

\begin{equation}
E^{\prime}=\left\{s^{\prime}=T\left(s^{\prime}, s\right) P_{E}(s) \quad \forall s \in E\right\}
\end{equation}

Now \(E^{\prime}\) is closer to \(E_\text{eq}\) than \(E\)

\begin{align}
    \|E' - E_\text{eq}\| &= \sum_{s'} |P'(s') - P_\text{eq}(s')| \nonumber \\ 
    &=\sum_{s'}\left|\sum_s T(s',s)\{P(s) - P_\text{eq}(s)\}\right| \nonumber \\
    &< \sum_{s,s'}T(s',s)|P(s) - P_\text{eq}(s)| \nonumber \\
    &= \sum_s |P(s) - P_\text{eq}(s)| \nonumber \\
    &= \|E - E_\text{eq}\|,
\end{align}
where the first equality follows from the definition of $E'$ and the definition of the fixed point, the inequality follows from the facts that $T\geq 0$ and that $P, P_\text{eq}$, being probabilities, cannot satisfy $P > P_\text{eq}$ everywhere, and then we make use of $\sum_sT(s',s) = 1$.

\begin{itemize}
  \item For a Markov process satisfying the above criteria, Perron's theorem  for positive matrices guarantees that $T$, viewed as a matrix in state space, has eigenvalues
\end{itemize}

\begin{equation}
\lambda_{1}=1>\lambda_{2} \geqslant \lambda_{3} \geqslant \ldots .
\end{equation}

\begin{itemize}
  \item The typical behaviour of a Markov process is best exemplified by considering an ensemble of many instances of the process starting from different initial samples and evolving in concert (all processes have the same \(T\) )
  \begin{itemize}
  \item After a few steps, the samples that are drawn will be far from the equilibrium distribution (they will represent samples from the initial distribution). As more steps are taken the distribution will thermalise to the target distribution
  \item Once the distribution has thermalised, new samples drawn will be representative of the target distribution. However evolving many Markov processes is expensive
  \item For an individual Markov process after thermalisation, sequences of samples will be representative of the target distribution. However these sequences will have autocorrelation (in most cases): nearby samples in the sequence will be closely related
  \end{itemize}
\end{itemize}

For a sequence of states \(\left\{s_{1}, \ldots  s_{N}\right\}\) and a quantity \(A(s)\) the integrated autocorrelation time is approximated by

\begin{equation}
\tau_{\text {int}, A}=\frac{1}{2}+\sum_{j=1}^{N} \Gamma_{A}(J)
\end{equation}

where

\begin{equation}
\Gamma_{A}(j)=C_{A}(j) / C_{A}(0)
\end{equation}

and the autocorrelation function is

\begin{equation}
C_{A}(j)=\left\langle\left(A_{i}-\langle A\rangle\right)\left(A_{i+j}-\langle A\rangle\right)\right\rangle
\end{equation}
with $A_i \equiv A(s_i)$ and $\langle A \rangle = \frac{1}{N}\sum_i A_i$. NB: this is observable dependent, we can define the exponential autocorrelation time as $\tau_\text{exp} = \operatorname{sup}_A \tau_{\text{int},A}$.

For correlated samples, the variances of samples is

\begin{equation}
\sigma_{A, \text{corr}}^{2}=\frac{1}{N^{2}} \sum_{i, j} C_{A}(|i-j|) \simeq 2 \tau_{\text{int}, A} \sigma_{A}^{2}
\end{equation}

where

\begin{equation}
\sigma_{A}^{2}=\left\langle\frac{1}{N^{2}}\left(\sum_{i=1}^{N}\left(A_{i}-\langle A\rangle\right)\right)^{2}\right\rangle
\end{equation}
and thus the effective number of independent samples for quantity $A$ is $N_\text{indep} = N/(2 \tau_{\text{int},A})$.
\begin{itemize}
    \item These autocorrelations must be accounted for in reliable statistical analysis. Alternatively, samples can be retained only after each $2 \tau_{\text{int},A}$ updates
\end{itemize}

\begin{figure}
    \centering
    \redrawcompare{fig-notes-079.png}{fig-redraw-079.pdf}
    \caption{A depiction of a Markov chain showing the initial thermalisation time followed by fluctuations about the equilibrium, with approximately independent samples being gathered after each autocorrelation time.}
    \label{fig:markov_chain}
\end{figure}

\section{Heatbath Algorithm}
\begin{itemize}
    \item Since $P_\text{eq} = \frac{1}{Z}e^{-S_\text{eff}(s)}$ is known up to the normalisation, choosing the transition probability
\end{itemize}
\begin{equation}
    T_{\text{HB}}(s', s) = \frac{1}{Z} e^{-S_\text{eff}(s')}
\end{equation}

satisfies detailed balance trivially:
\begin{align}
    T_{\text{HB}}(s', s)P_\text{eq}(s) &= \frac{1}{Z^2} e^{-S_\text{eff}(s')}e^{-S_\text{eff}(s)} \nonumber \\
    &=T_{\text{HB}}(s,s') P_\text{eq}(s')
\end{align}

\begin{itemize}
  \item However given that $s$ represents an entire field configuration over continuous degrees of freedom in our context there is no way to efficiently use such a global heatbath efficiently
\end{itemize}

\section{Metropolis Algorithm}
\begin{itemize}
    \item An algorithm satisfying detailed balance can be constructed as follows from an arbitrary initial state $s$
\end{itemize}
\begin{enumerate}
    \item Choose an updated state $s'$ according to an efficiently implementable transition probability $T_0(s', s)$
    \item Accept or reject the proposed sample according to
\end{enumerate}
\begin{equation}
    T_\text{ac}(s',s) = \operatorname{min}\left(1, \frac{T_0(s,s')P_\text{eq}(s')}{T_0(s',s)P_\text{eq}(s)}\right)
\end{equation}
\begin{itemize}
    \item The combined transition probability $T(s',s) = T_\text{ac}(s',s) T_0(s',s)$ satisfies detailed balance:
\end{itemize}
\begin{align}
T(s',s)P_\text{eq}(s) &= T_0(s',s) \operatorname{min}\left(1, \frac{T_0(s,s')P_\text{eq}(s')}{T_0(s',s)P_\text{eq}(s)}\right)P_\text{eq}(s) \nonumber \\
&= \operatorname{min}(T_0(s',s)P_\text{eq}(s),T_0(s,s')P_\text{eq}(s'))\nonumber \\
&= T_0(s,s')\operatorname{min}\left(\frac{T_0(s',s)P_\text{eq}(s)}{T_0(s,s')P_\text{eq}(s')},1\right)P_\text{eq}(s') \nonumber \\
&= T(s,s') P_\text{eq}(s')
\end{align}
\begin{itemize}
    \item A common choice is to have a symmetric $T_0(s,s') = T_0(s',s)$ such that
\end{itemize}
\begin{equation}
    T_\text{ac}(s', s) = \operatorname{min}(1-P_\text{eq}(s')/P_\text{eq}(s)) = \operatorname{min}(1, \exp{(-\Delta S)}),
\end{equation}
where $\Delta S \equiv S_\text{eff}(s') - S_\text{eff}(s)$.
\begin{itemize}
    \item For a field theory, an update $T_0$ can be to change a single degree of freedom in some specified way. For local actions this will change the action only locally so $T_\text{ac}$ is efficient to compute (cost $\approx \mathcal{O}(V)$).
    \begin{itemize}
        \item This can be repeated sweeping through each d.o.f. and then the acceptance step be performed, or acceptance can be performed immediately (but then the new configuration will be very correlated with the old one)
        \item For theories involving fermions $S_\text{eff}[\phi] \supset \operatorname{tr}\ln (D[\phi])$ is very nonlocal and so updating a single d.o.f. requires computing the entire $S_\text{eff}$, so the cost of the algorithm is $\mathcal{O}(V^2)$ - this is impractical for e.g. QCD
    \end{itemize}
\end{itemize}

\section{Langevin algorithm}
{ TODO}

\section{Pseudofermions}

\begin{itemize}
  \item For a complex scalar field
\end{itemize}

\begin{equation}
    \int \mathcal{D}\phi \mathcal{D}\phi^\dagger \exp\left[-\int \phi^\dagger A \phi\right] = \frac{c}{\det[A]} = c \det[A^{-1}],
\end{equation}
for $A$ a positive definite Hermitian matrix and $c = \pi^{\text{dim }A}$ constant (note in particular that this assumes $A$ has no zero eigenvalues).
\begin{itemize}
    \item Thus scalar ``pseudofermion'' fields can be used to represent the fermion determinant, so for fermions + gauge fields
\end{itemize}
\begin{align}
    Z &= \int \mathcal{D}U\mathcal{D}\bar{\psi}\mathcal{D}\psi e^{-S[U] - \int \bar{\psi}D[U]\psi} = \int \mathcal{D}U \det[D[U]]e^{-S[U]} \nonumber \\
    &= \int \mathcal{D}U \mathcal{D}\phi \mathcal{D}\phi^\dagger \exp \left\{ -S[U] - \int \phi^\dagger D^{-1}[U]\phi \right\}.
\end{align}

\begin{itemize}
    \item This assumes that $D$ is Hermitian positive definite.
\end{itemize}
\begin{enumerate}
    \item This is the case if there are an even number of fermion flavours with degenerate masses (e.g. $m_u = m_d = m_l)$
    \begin{equation}
        \det D_u \det D_d = (\det D_l)^2 = \det D_l \det D_l^\dagger = \det(D_l D_l^\dagger) > 0
    \end{equation}
    \item It also holds for a single flavour for a Dirac operator obeying the GW relation since eigenvalues come in complex conjugate pairs (also staggered fermions)
    \item It does not hold for a single flavour of Wilson fermions, since $\det D_W \in \mathbb{R}$ but not guaranteed to be positive. However for large enough mass and lattice volume, negative eigenvalues are exponentially suppressed and it is typical to use for e.g. strange quarks $\det D_s \approx [\det D_s^\dagger D_s]^{1/2}$
    \item In the presence of fermion chemical potential or the QCD $\theta$-term the determinant is nonpositive and importance sampling fails
\end{enumerate}

\begin{itemize}
  \item Since \(D^{-1}[U]\) is very non-local, the Metropolis algorithm will be very inefficient
\end{itemize}


\section{Hybrid Monte-Carlo Algorithm}
\begin{itemize}
    \item HMC introduces auxiliary canonical momenta conjugate to the field degrees of freedom and uses Hamilton's equations to evolve the system according to molecular dynamics
    \begin{itemize}
        \item This algorithm was developed for QCD simulations (Duane et al. 1987) but now sees much more use in statistics under the name of Hamiltonian Monte Carlo
    \end{itemize}
    \item Given a Hamiltonian $H(p,q) = \frac{1}{2}p^2 + S(q)$, the classical equations of motion are
    \begin{equation}
        \dot{p} = - \frac{\partial H}{\partial q} = - \frac{\partial S}{\partial q}, \qquad \dot{q} = \frac{\partial H}{\partial p} = p,
    \end{equation}
    where $q$ and $p$ are a set of variables and conjugate momenta. By integrating these equations, a trajectory in phase space $(q,p)$ is mapped out along which $H$ is constant.
    \item HMC uses these equations to propose new ``$q$''s that will have constant action and the update probability satisfies ergodicity and detailed balance provided the integration along the HMC trajectory is symplectic (area preserving) and reversible
    \item Specifying to a theory with $N_f = 2$ fermions and gauge fields, we can write a modified partition function (just multiplying by a constant)
\end{itemize}
\begin{equation}
    Z' = \int \mathcal{D}U \mathcal{D}P \mathcal{D}\phi \mathcal{D}\phi^\dagger \exp\left\{-S[U] - \sum_{n.m}\phi_n^\dagger (D^\dagger D[U])^{-1}\phi_m - \operatorname{Tr}\sum_{\mu,n} P_\mu^2(n) \right\},
\end{equation}
where the $P_\mu(n) \in g$ are elements of the Lie algebra $g$ and are conjugate to $Q_\mu(n)$ where $U_\mu(n) = \exp(i Q_\mu(n))$; note that the pseudofermion fields have spin and colour components.

\begin{itemize}
    \item Since correlation functions will never contain $P_\mu$'s they are unchanged
    \item The $P_\mu^a(n)$ are independent variables and are distributed as Gaussians of unit variance
\end{itemize}
\begin{itemize}
    \item The state of this extended theory is described by $s = \{ P_\mu(n), \phi(n), \phi^\dagger(n), U_\mu(n)\}$ and HMC provides updates between states.
\end{itemize}
This can be done in $4$ steps:
\begin{enumerate}
    \item Refresh pseudofermions: defining $\eta = D^{\dagger^{-1}}\phi$, $\eta^\dagger = \phi^\dagger D^{-1}$, the PF action is purely Gaussian for each component of $\eta_\alpha^a(n)$, so the PF update is done by generating Gaussian noise for $\eta$ and computing $\phi = D^\dagger \eta$
    \item Refresh momenta: also Gaussianly distributed for each $P_\mu^a(n)$
    \item Update momenta and gauge fields using Hamilton's equations: consider the canonical pairs $\{Q_\mu^a(n), P_\mu^a(n)\}$ where $U_\mu(n) = \exp(i Q_\mu(n))$ and the Hamiltonian
    \begin{equation}
        H(P,Q) = \sum_{\mu,n} \operatorname{Tr}(P_\mu^2(n)) + \phi^\dagger (D^\dagger D[Q])^{-1}\phi + S[Q],
    \end{equation}
    then $H$-conserving trajectories are found by integrating the fictitious time evolution of
    \begin{equation}
        \dot{P}_\mu(n) = -\frac{\partial H}{\partial Q_\mu(n)}, \qquad \dot{Q}_\mu(n) = \frac{\partial H}{\partial P_\mu(n)} = P_\mu(n)
    \end{equation}
    using a symplectic, reversible integrator (e.g. leapfrog integrator). This deterministically defines $T_{MD}(\{P',Q'\},\{P,q\})$.
    \item Accept-reject step: Using
    \begin{equation}
        T_\text{ac}(\{P',Q'\},\{P,Q\}) = \operatorname{min}(1,\exp(-H(P',Q')+H(P,Q)))
    \end{equation}
    Finally, convert back to $U_\mu' = \exp(i Q_\mu')$
\end{enumerate}


\begin{itemize}
  \item The transition probability
\end{itemize}

\begin{equation}
T(Q',Q) = \int \mathcal{D}P \mathcal{D}P'T_\text{ac}(\{P',Q'\},\{P,Q\})e^{-\sum_{\mu, n} P_\mu^2(n)}T_\text{MD}(\{P',Q'\},\{P,Q\})
\end{equation}

satisfies detailed balance (proof in G\& L).
\begin{itemize}
    \item A symplectic integrator such as leapfrog is reversible and area preserving as required for Hamiltonian system. For each integration step of length $\varepsilon$,
    \begin{enumerate}
        \item $P(\varepsilon/2) = P(0) - \varepsilon/2 \frac{\partial H}{\partial Q}|_{Q(0)}$
        \item $Q(\varepsilon) = Q(0) + \varepsilon \frac{\partial H}{\partial P}|_{P(\varepsilon/2)}$
        \item $P(\varepsilon) = P(\varepsilon/2) - \frac{\varepsilon}{2}\frac{\partial H}{\partial Q}|_{Q(\varepsilon)}$
    \end{enumerate}

\begin{figure}
    \centering
   \redrawcompare{fig-notes-021.jpg}{fig-redraw-021.pdf}
    \caption{A schematic representation of leapfrog integration; first, $P$ takes step $a$ with a step size $\varepsilon/2$, $Q$ leapfrogs it to $Q(\varepsilon)$ using $P(\varepsilon/2)$, and $P$ then steps to $\varepsilon$ using the value of $Q$ evaluated there.}
    \label{fig:leapfrog}
\end{figure}

    \item The transformation for each step is triangular so the Jacobian is trivial, ($dP dQ = dP' dQ'$)
    \item It is also easy to see that this integration is reversible
    \item Numerical errors in the integration mean that $\Delta H \neq 0$ along the integrated trajectory (leapfrog has $\mathcal{O}(\varepsilon^2)$ errors, but higher-order integrators exist). Without such errors, there would be no need for the accept-reject step. These errors also accumulate as the trajectory length grows, so in practice one balances how far in configuration space one moves with the acceptance
    \item The forces on momenta are complicated to compute


\begin{equation}
F_{\mu}=-\frac{\partial H}{\partial Q_{\mu}}=-\frac{\partial}{\partial Q_{\mu}}\left[S[Q]+\phi\left(D^{\dagger} D\right)^{-1} \phi\right]
\end{equation}
and the latter term involves evaluating $\chi = (D^\dagger D)^{-1}\phi$ which is equivalent to solving $D^\dagger D \chi = \phi$; this is the most computationally demanding part of HMC

    \item The algorithm scales as $\mathcal{O}(V^{5/4})$ using a second order symplectic integrator as acceptance drops as $V$ increases, requiring smaller integration time steps
    \item Improving (or replacing) HMC and its variants for single quark flavours is an active area of research
\end{itemize}

\section{Linear System Solvers}
\begin{itemize}
    \item In calculating correlation functions with fermion fields after integrating out the fermion fields we are left with products of $D^{-1}[U]$ acting on e.g. a spin-colour vector $\rho$ which can be computed by solving
\begin{equation}
    D(x,y) G(y) = \rho(x) \quad \Longrightarrow \quad G(y) = D^{-1}(x,y)\rho(x)
\end{equation}
\item We also saw that for HMC we need to solve
\begin{equation}
    (D^\dagger D)(x,y) \chi(y) = \phi(x)
\end{equation}
   \item $D$ and $D^\dagger D$ are large, $12 V \times 12 V$, sparse ($\mathcal{O}(30)$ nonzero entries per row) matrices so this is equivalent to solving the linear system
\end{itemize}
\begin{equation}
    M x = b
\end{equation}

\begin{itemize}
    \item The most efficient direct solvers (e.g., LU factorisation) scale as $\mathcal{O}(N^3)$ for an $N \times N$ linear system whcih is impractical for $V \sim 10^8$
    \item Iterative methods  construct successively improving approximations to $x$: $x_1 \rightarrow x_2 \rightarrow \ldots \rightarrow x_N$ with $\|x_{i+1} - \hat{x} \| < \|x_{i} - \hat{x} \|$ where $\hat{x}$ is the exact solution
    \item Krylov solvers are a class of iterative methods that work well for the sparse systems under consideration and scale only as $\mathcal{O}(N)$
    \item The conjugate gradient algorithm (CG) works well for Hermitian systems (NB: $D$ is not Hermitian but $D^{-1} = D^\dagger (D D^\dagger)^{-1}$, so it suffices to work with $A = D D^\dagger$ which is Hermitian)
    \begin{itemize}
        \item CG proceeds via minimising the quadratic form
        \begin{equation}
            f(\vec{x}) = \frac{1}{2}\vec{x}^\dagger A \vec{x} - \vec{x}^\dagger\vec{b},
        \end{equation}
        since at $\nabla f = 0$ we have $A \vec{x} = \vec{b}$.
        \item The minimisation is performed iteratively in successively orthogonal directions. It does this by finding $\alpha_i$ such that $f(\vec{x}_{(i+1)}=\vec{x}_{(i)}+\alpha_i\vec{p}_{(i)})$ is minimal with $\vec{p}_{(i)}$ orthogonal to all previous $\vec{p}_{(j)}$ in the A-norm
        \begin{equation}
            \vec{p}_{(i)}^\dagger A \vec{p}_{(j)} \sim \delta_{ij}
        \end{equation}
        This sequentially builds up the $i$th Krylov space
        \begin{equation}
            K_{(i)} = \operatorname{span}\{\vec{p}_{(0)} = \vec{b} - A \vec{x}_{(0)}, \vec{p}_{(1)}, \ldots, \vec{p}_{(i)}\}
        \end{equation}
        \item The sequence $\vec{x}_{(0)} \rightarrow \vec{x}_{(1)} \rightarrow \ldots \rightarrow \vec{x}_{(N)} = \hat{x}$ arrives at the exact solution after $N = \operatorname{dim}(A)$ steps
        \item Typically however $|\vec{x}_{(i)} - \hat{x}|$ is $0$ to within machine precision much earlier
    \end{itemize}
\end{itemize}

CG Algorithm$(A, \vec{b}, \text{tol})$:
\begin{algorithmic}
    \State $\vec{x}_{(0)} \gets$ random initialization
    \State $\vec{r}_{(0)} \gets \vec{b} - A \vec{x}_{(0)}$
    \State $\vec{p}_{(0)} = \vec{r}_{(0)}$
    \For{$k=0, \ldots, N$}
    \State $\alpha_k \gets \vec{r}_{(k)}^\dagger \cdot \vec{r}_{(k)} / (\vec{p}_{(k)}^\dagger A \vec{p}_{(k)})$
    \State $\vec{x}_{(k+1)} \gets \vec{x}_{(k)} + \alpha_k \vec{p}_{(k)}$
    \State $\vec{r}_{(k+1)} \gets \vec{r}_{(k)} - \alpha_k A \vec{p}_{(k)}$
    \If{$\| \vec{r}_{(k+1)}\| < \text{tol} \| \vec{b}\|$} 
    \State stop and return $\vec{x}_{(k+1)}$
    \EndIf
    \State $b_k \gets \vec{r}_{(k+1)}^\dagger \cdot \vec{r}_{(k+1)}/\vec{r}_{(k)}^\dagger \cdot \vec{r}_{(k)}$
    \State $\vec{p}_{(k+1)} = \vec{r}_{(k+1)} + b_k \vec{p}_{(k)}$
    \EndFor
\end{algorithmic}
\begin{itemize}
    \item CG is implemented solely in terms of matrix-vector and vector-vector operations
    \item Typically the number of iterations required for convergence to a particular precision is set by the condition number of the matrix $A$
\end{itemize}
\begin{equation}
    C(A) \equiv \frac{\lambda_\text{max}}{\lambda_\text{min}} = \text{ratio of largest to smallest eigenvalue}
\end{equation}
\begin{itemize}
    \item Since the lowest eigenvalue of $D^\dagger D$ is set (up to gauge field fluctuations) by $m_q$ light quarks are more challenging numerically
    \item Other variants are appropriate for non-Hermitian cases and may have faster convergence (BiCG-stab)
    \item Matrices can be preconditioned in many ways (even-odd, ILU, SSGR, Schwartz, Multigrid) to accelerate convergence
    \item The low eigenvalues of the matrix can be ``deflated'' (removed), also accelerating convergence
\end{itemize}

\section{Other algorithmic tasks}
\begin{itemize}
    \item Overlap fermions require the matrix sign function which is most optimally approximated using Zolotarev rational functions
    \item It is also necessary to compute $\operatorname{Tr}(D^{-1}(x,x))$ in some cases: this uses algorithms similar to linear solvers
\end{itemize}
